\documentclass[exercice]{thedocument}%
\usepackage{tabularray}
\startdatakeys
    \source{26GENMATME1}
    \BO{2020}
    \cycle{4}
    \competencesDuSocle{CA,MO}
    \connaissancesRequises{B1b*}
    \competencesTravaillees{B1b*,B3e*}
\enddatakeys
\begin{document}%
Le tableau ci-dessous présente le nombre de médailles obtenues
par 9 pays lors des Jeux Paralympiques de Paris 2024.
\begin{center}
    \begin{tblr}{%
            colspec={X[2,l] X[1,c] X[1,c] X[1,c] X[1,c]},
            hlines, vlines
        }
        Pays & Or & Argent & Bronze & Total \\
        Chine & 94 & 76 & 50 & 220 \\
        Grande-Bretagne & 49 & 44 & 31 & 124 \\
        États-Unis & 36 & 42 & 27 & 105 \\
        Pays-Bas & 27 & 17 & 12 & 56 \\
        Brésil & 25 & 26 & 38 & 89 \\
        Italie & 24 & 15 & 32 & 71 \\
        Ukraine & 22 & 28 & 32 & 82 \\
        France & 19 & 28 & 28 & 75 \\
        Australie & ? & 17 & 28 & 63
   \end{tblr}
\end{center}
\begin{enumerate}
    \item Comben de médailles d'or ont été obtenues par les Pays-Bas~?
    \item Calculer le nombre de médailles d'or obtenues par l'Australie.
    \item Montrer que l'affirmation suivante est vraie.\par\noindent
    \guillemets{Plus de 20\% des médailles obtenues par la Grande-Bretagne
    sont en bronze.} 
    \item \begin{enumerate}
        \item Déterminer la médiane de la série des nombres de médailles
        obtenues par ces 9 pays (colonne \guillemets{Total} du tableau).
        \par\noindent\textbf{Détailler la réponse en précisant la démarche}
        \item Intepréter ce résultat dans le contexte de l'exercice.
    \end{enumerate}
    \item Aux Jeux Paralympiques de Tokyo en 2021, le Brésil avait
    obtenu 20 médailles d'argent. Pour ceux de Paris en 2024, il a
    obtenu 26 médailles d'argent.\par\noindent
    Quel est le pourcentage d'augmentation du nombre de médailles
    d'argent obtenues par le Brésil entre 2021 et 2024~?
\end{enumerate}
\showthesource
\end{document}